\documentclass{beamer}

\usepackage[utf8]{inputenc}
\usepackage[english]{babel}
\usepackage{textpos}
\usepackage{graphicx}
\usepackage{textcomp}
\usepackage{animate}
\usepackage{comment}
%\usepackage{multicol}
%\setlength{\columnsep}{5cm}
\usepackage{amsmath}
\usepackage{amsfonts}
\usepackage{verbatim}
\usepackage{varwidth}
\usepackage{caption}
\usepackage{dcolumn}
\usepackage{amssymb}
\usepackage{algorithm}
\usepackage{algorithmic}
%\usepackage{physics}
\usepackage{hyperref}
\usepackage{scalerel,stackengine}
%\usepackage{fancyvrb}


\newcommand{\code}[1]{\texttt{#1}}
\newcommand{\shp}{$>>>$ }
\newcommand{\tab}{\hspace*{22pt}}

\usetheme{Madrid}

\title{Chapter 9: Exceptions}
\author{EEC 3150}
\date{}

\institute{Trevecca Nazarene University}
 


%%%%%%%%% BEGIN DOCUMENT %%%%%%%%%
%%%%%%%%% BEGIN DOCUMENT %%%%%%%%%
%%%%%%%%% BEGIN DOCUMENT %%%%%%%%%

\begin{document}

% get title page 
\frame{\titlepage}
 
%%%%%%%%%%
%%%%%%%%%%
\begin{frame}	
	\frametitle{Outline}
	\tableofcontents	
\end{frame}


%%%%%%%%%%
%%%%%%%%%%
\section{Handling Exceptions}

\begin{frame}	
	\center{\Huge{Handling Exceptions}}
\end{frame}

\begin{frame}	
	\frametitle{Handling Exceptions}
	
	\begin{alertblock}{Defn: exception}
		\textbf{Exceptions} are errors in a Python script detected by the interpreter
		
		\vspace{10pt}
		They come in two main classes:
		\begin{itemize}
			\item \emph{Syntax errors} occur when you write code that violates the rules of Python's syntax; detected by the interpreter before execution
			\begin{itemize}
				\item Missing a parenthesis or quote
				\item Misspelled keyword
				\item Using a variable before it's defined
			\end{itemize}
			\item \emph{Runtime errors} occur during the execution of a script; cannot be detected by the interpreter before execution
			\begin{itemize}
				\item Dividing by zero
				\item Out of bounds indexing of a list
				\item Trying to read a file that doesn't exist
			\end{itemize}
		\end{itemize}
		
		\vspace{5pt}
		Unless handled properly, exceptions always halt your program's execution. How can we avoid this?
	\end{alertblock}
	
\end{frame}

\begin{frame}	
	\frametitle{Handling Exceptions}
	
	\begin{alertblock}{Defn: exception handling}
		\emph{Exception handling} (or error handling) is the process of detecting and dealing with errors in programs
	\end{alertblock}
	
	\begin{alertblock}{Defn: the \code{try except} block}
		\minipage{0.6\textwidth}
			\begin{itemize}
				\item We can use a \code{try except} block to perform error handling in Python
				\item This will give us a plan B in case our code throws an error
				\item It's similar to an \code{if else} block
			\end{itemize}
			
			\vspace{6pt}
			\code{try: \\
					\tab \emph{code that might throw an error} \\
					except: \\
					\tab \emph{code to run in case of error}
			}
		\endminipage\hfill
		\minipage{0.4\textwidth}
			\begin{itemize}
				\item The code in the \code{try} portion is run first
				\item If no error, then the block is finished
				\item If there is an error, then the \code{except} portion is run
			\end{itemize}
		\endminipage\hfill
	\end{alertblock}
	
\end{frame}

\begin{frame}	
	\frametitle{Handling Exceptions}
	
	\begin{block}{\code{try except} example}
		\minipage{0.6\textwidth}
			Handling indexing a list out of bounds
			
			\vspace{10pt}
			\code{try: \\
					\tab \emph{code that might throw an error} \\
					except: \\
					\tab \emph{code to run in case of error}
			}
		\endminipage\hfill
		\minipage{0.4\textwidth}
			\begin{itemize}
				\item The code in the \code{try} portion is run first
				\item If no error, then the block is finished
				\item If there is an error, then the \code{except} portion is run
			\end{itemize}
		\endminipage\hfill
	\end{block}
	
\end{frame}






\end{document}




